\clearpage
\begin{explanationbox}
\textbf{第一步，先只看两个熟悉的直角三角形。}
普通点 $P$ 在斜线段 $AB$ 内；$H,K$ 是两端向过 $P$ 的水平直线所作垂足。
\begin{center}\begin{tikzpicture}[font=\small,>=stealth]
\draw[rounded corners=3pt,gray!45] (0,0) rectangle (14.2,12);
\node at (7.1,11.45) {先看初中几何：同一条斜线上的两个直角三角形相似};
\begin{scope}[shift={(1.2,4.7)},scale=1.4]

\fill[blue!7] (2,0)--(4,0)--(4,3.464101615)--cycle;
\fill[orange!10] (2,0)--(1.333333333,0)--(1.333333333,-1.154700538)--cycle;

\draw[gray,dashed] (4,3.464101615)--(4,0);
\draw[gray,dashed] (1.333333333,-1.154700538)--(1.333333333,0);
\draw[blue!75!black,very thick] (2,0)--(4,3.464101615) node[midway,above left=5pt] {$r$};
\draw[orange!85!black,very thick] (2,0)--(1.333333333,-1.154700538) ;\node[orange!85!black] at (2.35,-1.12) {$s$};
\draw[blue!65!black,thick] (2,0)--(4,0) node[midway,below=7pt] {$u$};
\draw[orange!85!black,thick] (1.333333333,0)--(2,0) node[midway,above=6pt] {$v$};
\draw[gray!80] (3.8,0)--(3.8,.2)--(4,.2);
\draw[gray!80] (1.513333333,0)--(1.513333333,-.18)--(1.333333333,-.18);
\draw[violet] (2.48,0) arc[start angle=0,end angle=60,radius=.48];
\draw[violet] (1.52,0) arc[start angle=180,end angle=240,radius=.48];
\node[violet] at (2.76,.28) {$\alpha$};
\node[violet] at (2.18,-.58) {$\alpha$};
\fill (4,3.464101615) circle(2pt) node[above=8pt] {$A$};
\fill (1.333333333,-1.154700538) circle(2pt) node[below left=6pt] {$B$};
\fill (4,0) circle(1.5pt) node[below right=5pt] {$H$};
\fill (1.333333333,0) circle(1.5pt) node[above=14pt] {$K$};
\fill (2,0) circle(2pt) node[below right=5pt] {$P$};
\end{scope}

\node[anchor=west,align=left] at (8.5,7.1)
 {$A,P,B$ 共线\\[9pt]$H,P,K$ 共线\\[9pt]$AH\parallel BK$\\[13pt]
 $PA=r,\quad PB=s$\\[10pt]
 $PH=u,\quad PK=v$\\[13pt]
 对应锐角同为 $\alpha$\\[9pt]
 两个直角三角形相似};
\node at (7.1,1.55) {$\triangle PAH\sim\triangle PBK$：比较同类边，不需要先学新曲线。};
\node at (7.1,.6) {$\displaystyle \frac{PH}{PA}=\frac{PK}{PB}\quad\Longrightarrow\quad\frac ur=\frac vs$};
\end{tikzpicture}
\end{center}
两个直角相等，$P$ 处两个锐角又是对顶角，所以 $\triangle PAH\sim\triangle PBK$。
对应的水平边与斜边之比相等：$u/r=v/s$。
\textbf{先看懂这个比；它来自相似，不需要背新的公式。}
\end{explanationbox}
\clearpage
\begin{explanationbox}
\textbf{第二步，把同一结构放回本题。}图形和尺度保持不变；把普通点 $P$ 对应为焦点 $F$，再补上曲线与准线。
\begin{center}\begin{tikzpicture}[font=\small,>=stealth]
\draw[rounded corners=3pt,gray!45] (0,0) rectangle (14.2,12);
\node at (7.1,11.45) {再认对应：相似关系保留，只补上抛物线的等距定义};
\begin{scope}[shift={(1.2,4.7)},scale=1.4]

\fill[blue!7] (2,0)--(4,0)--(4,3.464101615)--cycle;
\fill[orange!10] (2,0)--(1.333333333,0)--(1.333333333,-1.154700538)--cycle;

\draw[teal!75!black,thick,samples=120,domain=-1.9:3.9,variable=\t] plot ({1+\t*\t/4},{\t});
\draw[violet,densely dashed] (0,-1.9)--(0,3.95);
\node[violet,rotate=90] at (-.5,1.5) {准线};
\draw[blue!65!black,dashed] (0,3.464101615)--(4,3.464101615);
\draw[orange!85!black,dashed] (0,-1.154700538)--(1.333333333,-1.154700538);
\fill (0,3.464101615) circle(1.5pt) node[left=5pt] {$M$};
\fill (0,-1.154700538) circle(1.5pt) node[left=5pt] {$N$};
\fill (0,0) circle(1.5pt) node[left=5pt] {$Q$};
\draw[gray,dashed] (0,0)--(1.333333333,0);
\node[above=6pt] at (.6,0) {$p$};

\draw[gray,dashed] (4,3.464101615)--(4,0);
\draw[gray,dashed] (1.333333333,-1.154700538)--(1.333333333,0);
\draw[blue!75!black,very thick] (2,0)--(4,3.464101615) node[midway,above left=5pt] {$r$};
\draw[orange!85!black,very thick] (2,0)--(1.333333333,-1.154700538) ;\node[orange!85!black] at (2.35,-1.12) {$s$};
\draw[blue!65!black,thick] (2,0)--(4,0) node[midway,below=7pt] {$u$};
\draw[orange!85!black,thick] (1.333333333,0)--(2,0) node[midway,above=6pt] {$v$};
\draw[gray!80] (3.8,0)--(3.8,.2)--(4,.2);
\draw[gray!80] (1.513333333,0)--(1.513333333,-.18)--(1.333333333,-.18);
\draw[violet] (2.48,0) arc[start angle=0,end angle=60,radius=.48];
\draw[violet] (1.52,0) arc[start angle=180,end angle=240,radius=.48];
\node[violet] at (2.76,.28) {$\alpha$};
\node[violet] at (2.18,-.58) {$\alpha$};
\fill (4,3.464101615) circle(2pt) node[above=8pt] {$A$};
\fill (1.333333333,-1.154700538) circle(2pt) node[below left=6pt] {$B$};
\fill (4,0) circle(1.5pt) node[below right=5pt] {$H$};
\fill (1.333333333,0) circle(1.5pt) node[above=14pt] {$K$};
\fill (2,0) circle(2pt) node[below right=5pt] {$F$};
\end{scope}

\node[anchor=west,align=left] at (8.5,7.1)
 {普通点 $P\longrightarrow$ 焦点 $F$\\[13pt]
 焦准距 $FQ=p$\\[13pt]
 $\textcolor{blue!75!black}{r=AF=AM=p+u}$\\[10pt]
 $\textcolor{orange!85!black}{s=BF=BN=p-v}$\\[13pt]
 相似比仍是 $u/r=v/s$};
\node at (7.1,1.55) {几何给出比值相等，抛物线定义把投影换成 $r-p$ 与 $p-s$。};
\node at (7.1,.6) {$\displaystyle \frac{r-p}{r}=\frac{p-s}{s}\quad\Longrightarrow\quad\frac1r+\frac1s=\frac2p$};
\end{tikzpicture}
\end{center}
准线在左边，焦准距是 $p$。由等距定义，蓝段 $r=p+u$，橙段 $s=p-v$。
\[
 \frac ur=\frac vs
 \ \Longrightarrow\ \frac{r-p}{r}=\frac{p-s}{s}
 \ \Longrightarrow\ 1-\frac pr=\frac ps-1.
\]
于是 $p/r+p/s=2$。\textbf{倒数来自“水平边除以斜边”，再用定义把水平边换掉。}
\end{explanationbox}
\clearpage
\begin{explanationbox}
\textbf{第三步，转动焦点弦，看倒数怎样互相补偿。}这里统一按上下标号：上端为 $A$，下端为 $B$。
\begin{center}\begin{tikzpicture}[font=\small,>=stealth]
\draw[rounded corners=3pt,gray!45] (0,0) rectangle (14.2,12.5);
\node at (7.1,11.95) {两段长度在变，倒数一增一减：双色条的总长度始终相同};
\node at (2.55,10.95) {$\theta=60^\circ$};
\begin{scope}[shift={(1.9,6.6)},scale=0.65]
\draw[->,gray] (-1.2,0)--(4.3,0) node[above] {$x$};
\draw[teal!75!black,thick,samples=120,domain=-3.85:3.85,variable=\t] plot ({\t*\t/4},{\t});
\draw[blue!75!black,very thick] (1,0)--(3.000000000000001,3.4641016151377553);
\draw[orange!85!black,very thick] (1,0)--(0.33333333333333326,-1.1547005383792515);
\fill (3.000000000000001,3.4641016151377553) circle(2pt) node[above=8pt] {$A$};
\fill (0.33333333333333326,-1.1547005383792515) circle(2pt) node[below left=6pt] {$B$};
\fill (1,0) circle(2pt) node[below right=5pt] {$F$};
\end{scope}
\fill[blue!28] (1.2,2.15) rectangle (2.0375,2.55);
\fill[orange!35] (2.0375,2.15) rectangle (4.55,2.55);
\draw[gray] (1.2,2.15) rectangle (4.55,2.55);
\node at (1.61875,1.65) {$\textcolor{blue!75!black}{1/4}$};
\node at (3.29375,1.65) {$\textcolor{orange!85!black}{3/4}$};
\node at (2.55,3.25) {$1/AF+1/BF=1$};
\node at (6.95,10.95) {$\theta=90^\circ$};
\begin{scope}[shift={(6.3,6.6)},scale=0.65]
\draw[->,gray] (-1.2,0)--(4.3,0) node[above] {$x$};
\draw[teal!75!black,thick,samples=120,domain=-3.85:3.85,variable=\t] plot ({\t*\t/4},{\t});
\draw[blue!75!black,very thick] (1,0)--(1.0000000000000002,2.0000000000000004);
\draw[orange!85!black,very thick] (1,0)--(0.9999999999999999,-2.0);
\fill (1.0000000000000002,2.0000000000000004) circle(2pt) node[above=8pt] {$A$};
\fill (0.9999999999999999,-2.0) circle(2pt) node[below left=6pt] {$B$};
\fill (1,0) circle(2pt) node[below right=5pt] {$F$};
\end{scope}
\fill[blue!28] (5.6,2.15) rectangle (7.2749999999999995,2.55);
\fill[orange!35] (7.2749999999999995,2.15) rectangle (8.95,2.55);
\draw[gray] (5.6,2.15) rectangle (8.95,2.55);
\node at (6.4375,1.65) {$\textcolor{blue!75!black}{1/2}$};
\node at (8.112499999999999,1.65) {$\textcolor{orange!85!black}{1/2}$};
\node at (6.95,3.25) {$1/AF+1/BF=1$};
\node at (11.35,10.95) {$\theta=120^\circ$};
\begin{scope}[shift={(10.7,6.6)},scale=0.65]
\draw[->,gray] (-1.2,0)--(4.3,0) node[above] {$x$};
\draw[teal!75!black,thick,samples=120,domain=-3.85:3.85,variable=\t] plot ({\t*\t/4},{\t});
\draw[blue!75!black,very thick] (1,0)--(0.3333333333333336,1.1547005383792517);
\draw[orange!85!black,very thick] (1,0)--(2.9999999999999982,-3.4641016151377535);
\fill (0.3333333333333336,1.1547005383792517) circle(2pt) node[above=8pt] {$A$};
\fill (2.9999999999999982,-3.4641016151377535) circle(2pt) node[below left=6pt] {$B$};
\fill (1,0) circle(2pt) node[above right=5pt] {$F$};
\end{scope}
\fill[blue!28] (10.0,2.15) rectangle (12.5125,2.55);
\fill[orange!35] (12.5125,2.15) rectangle (13.35,2.55);
\draw[gray] (10.0,2.15) rectangle (13.35,2.55);
\node at (11.25625,1.65) {$\textcolor{blue!75!black}{3/4}$};
\node at (12.931249999999999,1.65) {$\textcolor{orange!85!black}{1/4}$};
\node at (11.35,3.25) {$1/AF+1/BF=1$};

\node at (7.1,.75) {$y^2=4x$，$p=2$；蓝色表示 $1/AF$，橙色表示 $1/BF$，总和为 $2/p=1$。};
\end{tikzpicture}
\end{center}
图中蓝、橙两段通常不等，取倒数后也不同；但两块条形合起来始终为 $1$。
这三条弦用于观察和核验，下一节再给出任意焦点弦的证明。
\[
 \frac1r=\frac{1-\cos\theta}{p},\qquad
 \frac1s=\frac{1+\cos\theta}{p}.
\]
一增一减的部分相加抵消，留下固定的 $2/p$。通径只是其中 $r=s=p$ 的对称情形。
\end{explanationbox}
